\section{Telemetría y órdenes}
\label{sec:telemetria}

\subsection{Datos que muestra la app}

\begin{table}[H]
\centering\small
\begin{tabularx}{\textwidth}{@{}l L@{}}
\toprule
\textbf{Dato} & \textbf{Cómo se muestra} \\
\midrule
Conectividad & Conectado, señal débil o sin datos, con la latencia en milisegundos \\
Batería & Porcentaje y minutos estimados. Ámbar debajo del 25~\%, rojo debajo del 15~\% \\
Estado operativo & En base, listo, patrullando, en pausa, incidente, regresando, detenido o falla \\
Misión y progreso & Ruta activa, punto de control actual y porcentaje recorrido \\
Ubicación & Posición sobre el plano y nombre de la zona \\
Movimiento & Velocidad y rumbo \\
Alertas & Cantidad de alertas activas y la más grave \\
Última actualización & Hora del último dato y su antigüedad. Con más de 3~s sin datos, todo pasa a gris con la leyenda ``hace 12~s'' \\
\bottomrule
\end{tabularx}
\end{table}

Las tres variables propias del servicio son:

\begin{enumerate}
  \item \textbf{Puntos de control verificados} en la ronda actual, por ejemplo 3 de 8. Es lo que el jefe de seguridad quiere ver al día siguiente.
  \item \textbf{Personas detectadas en zonas restringidas} durante la ronda.
  \item \textbf{Puertas abiertas detectadas}, que de noche tendrían que estar todas cerradas.
\end{enumerate}

\subsection{Órdenes}

Cada orden pasa por cuatro estados posibles, y la app siempre muestra en cuál está: \textit{pendiente} mientras espera respuesta, \textit{confirmada} si el robot la aceptó, \textit{rechazada} si la rechazó con un motivo, y \textit{sin respuesta} si pasaron 5~segundos sin noticias.

\begin{nota}
Que una orden esté confirmada quiere decir que el robot la aceptó, no que la terminó. Que el robot efectivamente se detuvo o llegó a la base se ve después, en la telemetría. Y la detención de la app es un comando de software simulado: no reemplaza la parada de emergencia física del robot.
\end{nota}

\begin{table}[H]
\centering\small
\begin{tabularx}{\textwidth}{@{}l L l@{}}
\toprule
\textbf{Orden} & \textbf{Se habilita cuando} & \textbf{Rechazo típico} \\
\midrule
Iniciar misión & Hay datos recientes, el robot está en base o listo, la batería está en 30~\% o más y no hay un incidente abierto & Batería baja, robot ocupado \\
Pausar & Está patrullando & Estado inválido \\
Reanudar & Está en pausa, con 15~\% de batería o más y sin incidente abierto & Batería baja \\
Cancelar misión & Patrullando, en pausa, en incidente o detenido. Pide confirmación & Estado inválido \\
Regresar a punto seguro & En cualquier estado salvo en base o ya regresando. Pide confirmación & Camino bloqueado \\
Detención operativa & Siempre. Hay que mantener apretado el botón 1,5~s & Puede quedar sin respuesta \\
Reconocer alerta & Hay una alerta activa sin reconocer & Alerta desconocida \\
Resolver incidente & Hay un incidente abierto. El guardia elige la decisión y después la confirma & Alerta desconocida \\
\bottomrule
\end{tabularx}
\end{table}

Mientras una orden está pendiente no se puede mandar otra del mismo tipo. La excepción es la detención operativa: se puede volver a mandar siempre, aunque haya una pendiente. Sin datos recientes los botones siguen visibles y habilitados, porque el guardia tiene que poder intentar detener el robot, pero la app avisa que la orden puede no llegar. Solo iniciar misión se deshabilita, porque necesita datos al día.

Cancelar se permite también con el robot detenido. Si no, después de una detención la única salida sería mandarlo a la base.

\subsection{Contrato de mensajes propuesto}
\label{sec:contrato}

Este contrato es una propuesta nuestra para un gateway propio entre la app y el robot. No es el protocolo oficial de Unitree. El repositorio simulado lo usa de verdad: serializa cada orden a JSON y parsea la telemetría, las respuestas y las alertas que le devuelve el simulador.

Los ejemplos que siguen son mensajes reales del simulador, tomados de una ronda de prueba. Los \texttt{timestamp} viajan en ISO-8601 y en UTC, así que 03:02Z son las 00:02 en Buenos Aires. Están en \texttt{docs/contrato} y un test unitario los compara con lo que genera el simulador, así que no se desactualizan. Los decimales largos de posición y progreso salen así del simulador.

\textbf{Telemetría}

\noindent
\begin{minipage}[t]{0.49\textwidth}
\small Patrullando por el pasillo A
\lstinputlisting[basicstyle=\ttfamily\scriptsize]{../contrato/telemetria-patrullando.json}
\end{minipage}\hfill
\begin{minipage}[t]{0.49\textwidth}
\small Batería por debajo del 25~\%
\lstinputlisting[basicstyle=\ttfamily\scriptsize]{../contrato/telemetria-bateria-baja.json}
\end{minipage}

\textbf{Orden y respuestas}

El primer intento de arranque fue con 20~\% de batería y el robot lo rechazó. Después de cargar, el segundo se aceptó.

\noindent
\begin{minipage}[t]{0.49\textwidth}
\lstinputlisting{../contrato/orden-iniciar-mision.json}
\end{minipage}\hfill
\begin{minipage}[t]{0.49\textwidth}
\lstinputlisting{../contrato/ack-aceptada.json}
\lstinputlisting{../contrato/ack-rechazada-bateria-baja.json}
\end{minipage}

\textbf{Alertas y decisión del guardia}

Ante el obstáculo el guardia eligió esperar. La orden para resolver el incidente lleva el id de la alerta y la decisión.

\noindent
\begin{minipage}[t]{0.49\textwidth}
\lstinputlisting{../contrato/alerta-obstaculo.json}
\lstinputlisting{../contrato/orden-resolver-incidente.json}
\end{minipage}\hfill
\begin{minipage}[t]{0.49\textwidth}
\lstinputlisting{../contrato/alerta-persona-zona-restringida.json}
\end{minipage}

Si pasan 5~segundos sin un \texttt{command\_ack}, la app marca la orden como sin respuesta. No se reintenta sola: el guardia decide si la vuelve a mandar.
